\subsection*{Grace makes the saints capable of charity}
The Collect asks for good works as the fruit of grace that goes before and follows. The Epistle gives the corresponding act of prayer: Paul kneels before the Father and asks strength through the Spirit, Christ's indwelling through faith and rootedness in charity. What the Collect asks for continually is the very life for which the apostle asks divine power. The Secret then names a sacramental means of cleansing and asks mercy to perfect the people for participation. These requests make activity and dependence mutually intelligible.

Aquinas's exposition of Ephesians~3:14 reads knees as the body's acknowledgment of weakness and smallness, and perseverance as a gift. Augustine's exposition of Psalm~70:16--18 refuses to make God dispensable once he has shown the way: the teaching from youth and the prayer into old age express a life whose initiation and continuance require him. The Communion brings that latter prayer to those receiving the sacrament. Their capacity to act is a received good, and its exercise never makes mercy superfluous. \sourcecite{Aquinas, \textit{Super Ephesios} 3.4; Augustine, \textit{Enarrationes} 70, abridged NPNF ``Psalm LXXI,'' §§18--21.}

Chrysostom's prayerful account joins the recipients' faith and love to gifts beyond their prior asking. The same passage begins with the apostle's affliction for them. Strength therefore concerns a charity capable of common life amid difficulty: Ephesians~2 gives the two peoples access in one Spirit, and chapter~4 requires humility, patience and unity. In the Gospel, Christ's actual relief of distress and his correction of ambition give good works concrete direction. \sourcecite{Chrysostom, \textit{Homilies on Ephesians} VII; Eph.~2:18--22; 4:1--6; Luke~14:1--11.}

\subsection*{The giver of honor frees the guest from rivalry}
Luke distinguishes the healing from the seating instruction. Jesus cures and dismisses the man before he addresses the guests; the wedding is the parable's imagined setting. No seat is assigned to the man who was healed. The guests, however, are shown choosing precedence, and the parable gives the inviter the power to confer a higher place. Received honor remains a good while self-exaltation receives correction.

Cyril grounds humility in truthful self-knowledge: wealth passes, bodily weakness and mortality persist, and even a claim to honor can be relinquished. Francis de Sales's last-place echo uncovers the heart that performs lowliness to obtain greater honor. His refusal of false humility that obstructs charity belongs to the same argument: humility should receive gifts truthfully and use them. \sourcecite{Cyril, \textit{Commentary on Luke}, Sermon CII, pp.~476--479; Francis de Sales, \textit{Introduction à la vie dévote} III.5.}

Paul's glory has a different setting. The apostle's sufferings serve the hearers, and the doxology gives glory to God in the Church. The Communion likewise remembers God's justice. Together these texts free good from the demand that it certify one's precedence: the host can honor a guest, the apostle can serve through affliction, and the communicant can acknowledge a received righteousness. The Offertory's threatened singer still asks protection and the shame of destructive enemies. Humility does not make injury desirable or confer a right to harm. \sourcecite{Eph.~3:13,21; Ps.~70:16 and 39:14--15, appointed selections; Aquinas, \textit{Super Psalmo} 39 n.~7.}

\subsection*{The heart's indwelling belongs to the Church's building}
The Epistle places Christ's indwelling within comprehension shared with all the saints. The Gradual answers with Sion built before nations and kings; the Alleluia calls for a new canticle in response to divine wonders. The \textit{Liturgical Year} continuation explicitly joins inward strengthening, ecclesial unity, Sion and praise. The conjunction already has a liturgical exposition, while each image keeps its own scriptural subject. \sourcecite{\textit{Liturgical Year}, 1909 continuation, pp.~359,363.}

Theodoret begins the Gradual psalm with the restoration of the humbled city. Its neighbors' amazement, followed by continuing unbelief and attacks, shows why the universal promise exceeds that restoration. Augustine reads the building as the Church's present gathering before Christ's glorious judgement, and Bellarmine contrasts that glorious appearance with the humility of his first coming. The city and the inward life therefore open toward an assembled people whose future exceeds the present work. \sourcecite{Theodoret, PG~80, 1679B--1682B; Augustine, \textit{Enarrationes} 101, abridged NPNF ``Psalm CII,'' §§16--18; Bellarmine, p.~317, the Gentiles and built-Sion lemmas.}

Chrysostom reads Paul's dimensions as the immensity of Christ's love. Augustine and Aquinas unfold a further reading through the Cross: charity works broadly, perseveres, hopes for heaven and receives its hidden beginning in grace. Aquinas also explains the dimensions through divine perfections and distinguishes knowing God's presence from containing all that God is. Charity can be known without being possessed as a measurable quantity. Benedict XVI develops this into an account of contemplation in which mind and heart enter Christ's inexhaustible mystery through love. \sourcecite{Chrysostom VII; Augustine, Letter~140, 25.62--26.64; Aquinas, \textit{Super Ephesios} 3.5; Benedict XVI, General Audience, 14 January 2009.}

The Alleluia makes praise the response to this divine work. Theodoret's new song belongs to the Savior's new manner of worship and life; Bellarmine traces the wonders through Christ's saving acts. The whole Psalm~97 makes Israel's salvation visible before the nations and summons creation to welcome just judgement. The new canticle has that public breadth, not merely the freshness of an individual's sensation. \sourcecite{Theodoret, PG~80, 1657C--1658D; Bellarmine, p.~306; Ps.~97:1--9.}

\subsection*{Divine assistance embraces mind and body}
The man in the Gospel receives a bodily cure. The saints in the Epistle receive inward strength; the Secret asks cleansing through the present sacrifice and the Postcommunion renewal through heavenly mysteries with bodily assistance now and hereafter. The forms of help differ, yet bodily distress remains within the prayer for the whole person's good.

Cyril presents Sabbath mercy as revealing the law's divine purpose and holy action as belonging to spiritual rest from sin. His son/ox text differs from the Missal's ass/ox, and his appeal to parental affection depends on his own witness. Augustine distinguishes bodily recovery and inward healing when expounding the continuation of the Alleluia's verse. His Gospel allusion there concerns Nain in Luke~7, not the present Luke~14 healing. Each exposition joins Christ's power to the reform of life without treating physical disease as evidence of personal guilt. \sourcecite{Cyril, Sermon CI, pp.~471--475; Augustine, \textit{Enarrationes} 97, abridged NPNF ``Psalm XCVIII,'' §1.}

The sacramental pair has an actual history together: the Gelasian book III, section XII sets \latin{Munda nos} and \latin{Purifica} in their offering and closing positions, with different opening prayers. The Gregorian supplement groups the present three orations under Sunday XVII. The \textit{Liturgical Year} continuation develops their sequence toward Communion and renewal affecting the body in this life and the next. Eucharistic participation is not identified with the Gospel's immediate cure, and the future bodily good exceeds the preservation of present health. \sourcecite{Wilson, \textit{Gelasian}, p.~231; \textit{Gregorian}, p.~174; \textit{Liturgical Year}, pp.~370--371.}

\subsection*{Remembered mercy becomes renewed petition}
The Introit's all-day cry, the Offertory's repeated request for aid and the Communion's old-age plea persist alongside Paul's praise of abundant power. Their wider psalms remember help already received and still name enemies, weakness and threatened life. The Postcommunion continues prayer after reception. Thanksgiving and petition belong to the same dependence on the giver.

Augustine hears the Introit as Christ and his members praying through successive generations. In the Offertory psalm he hears suffering members call to their physician. His conversion reading of the enemy's backward turn concerns a clause that the chant omits. Theodoret emphasizes the thwarting of persecution; Aquinas retains both penitential and punitive shame under divine justice. The request remains a serious plea that destructive harm be defeated. \sourcecite{Augustine, \textit{Enarrationes} 85, abridged NPNF ``Psalm LXXXVI,'' §§1--3,5,7; 39, ``Psalm XL,'' §§21--25; Theodoret, PG~80, 1159A--1160D; Aquinas, \textit{Super Psalmo} 39 n.~7.}

The Communion's ages also have distinct received interpretations. Augustine develops conversion and perseverance, applying age to a person and to the Church. Theodoret reads the people's schooling through Moses and the oldness of the law toward testimony to the Church from the nations; some added clauses in the PG text are bracketed. Bellarmine hears David's personal reliance on God's power and continuing testimony. The chant's omitted next-generation purpose illuminates the whole psalm but does not extend its appointed wording. Remembrance of God's help prepares further praise and further asking, through the vulnerability the giver continues to sustain. \sourcecite{Augustine, \textit{Enarrationes} 70, abridged NPNF ``Psalm LXXI,'' §§18--21; Theodoret, PG~80, 1423B--1426C; Bellarmine, pp.~212--213.}
